%% -----------------------------------------------------------
%% Relatorio.tex
%% Relatório do Trabalho Prático — Simulador de Paginação
%% Disciplina: Sistemas Operacionais — Ciência da Computação — UEM
%%
%% Formatação: template SBC (sbc-template.sty).
%% Compilar com pdflatex (inputenc utf8) + bibtex.
%% -----------------------------------------------------------

\documentclass[12pt]{article}

\usepackage{sbc-template}

\usepackage{graphicx,url}

\usepackage[brazilian]{babel}
\usepackage[T1]{fontenc} % Codificação para fontes com suporte a caracteres acentuados
\usepackage[utf8]{inputenc} % Codificação UTF-8 (necessário apenas para pdflatex)
\usepackage{multicol}
\usepackage{multirow}
\usepackage{amsmath,amssymb,amsfonts}
\usepackage{textcomp}
\usepackage{verbatim}
\usepackage[table,xcdraw]{xcolor}
\usepackage[shortlabels]{enumitem}
\usepackage{comment} % pacote de comentarios
%% O cabeçalho do SBGames não se aplica a um relatório de disciplina.
%% Para reativá-lo, descomente as duas linhas marcadas com [SBGames]
%% e preencha sbgames_info.tex.
%\usepackage{sbgames} % [SBGames]

%previne hifenização
\tolerance=1
\emergencystretch=\maxdimen
\hyphenpenalty=10000
\hbadness=10000

\sloppy

%% listas mais compactas, para caber no limite de páginas
\setlist[itemize]{topsep=3pt,itemsep=1pt,parsep=0pt,leftmargin=*}
\setlist[enumerate]{topsep=3pt,itemsep=1pt,parsep=0pt,leftmargin=*}

%
\trilha{Informe a Trilha} %Preencha com: Artes \& Design | Computação | Cultura | Saúde | Indústria | Educação | Tutoriais | MAGICA | WGP-Jogos
\local{Cidade/UF} % Exemplo: Manaus/AM
\ano{20XX} % Exemplo: 2024
\edicao{Edição (Romano)} %Exemplo: XXIII % [SBGames]

\title{Simulador de Paginação de Memória Virtual:
planejamento do trabalho \\
\medskip
\small{
    \textit{Title: A Virtual Memory Paging Simulator: work plan}
}
}

\author{Pedro Luiz Pereira Zampar\inst{1}, Pedro Augusto dos Santos Watermann\inst{1}}

\address{Departamento de Informática -- Universidade Estadual de Maringá (UEM)\\
  Av. Colombo, 5790 -- 87.020-900 -- Maringá -- PR -- Brasil
  \email{ra145117@uem.br, ra145121@uem.br}
}


\begin{document}

\maketitle

\thispagestyle{plain}

\begin{abstract}
This document reports the plan for a program that simulates how the operating
system manages memory when it does not fit entirely in physical memory. Given a
page reference string and a number of frames, the simulator counts page faults
and prints the state of the frames at each reference. The same run is repeated
under FIFO, LRU, Optimal, Second Chance and LFU, so that the five policies can
be compared as the amount of memory varies. We describe the file layout, the
function signatures, the expected behaviour of the program, the tests we intend
to run --- including Belady's Anomaly --- and the division of the work.
\end{abstract}

\keywords{Operating Systems, Virtual Memory, Paging, Page Replacement,
Belady's Anomaly.}

\begin{resumo}
Este documento registra o planejamento de um programa que simula como o sistema
operacional gerencia a memória quando ela não cabe toda na memória física. A
partir de uma sequência de páginas acessadas e de um número de quadros, o
simulador conta as faltas de página e mostra o estado dos quadros a cada acesso.
A mesma execução é repetida com FIFO, LRU, Ótimo, Segunda Chance e LFU, de modo
a comparar as cinco políticas variando a quantidade de memória. Descrevemos a
organização dos arquivos, as assinaturas das funções, como o programa deve se
comportar, os testes que pretendemos fazer --- incluindo a Anomalia de Belády
--- e a divisão do trabalho.
\end{resumo}

\palavraschave{Sistemas Operacionais, Memória Virtual, Paginação, Substituição
de Páginas, Anomalia de Belády.}

% =====================================================================
\section{O que vamos construir}
\label{sec:intro}

O trabalho é um programa que simula como o sistema operacional gerencia a
memória quando ela não cabe toda na memória física. O programa recebe uma lista
de páginas que o processo acessou e a quantidade de quadros disponíveis. A
partir disso, ele conta quantas vezes a página pedida não estava na memória (as
faltas de página) e mostra o passo a passo do que aconteceu.

A ideia é rodar essa mesma simulação com cinco algoritmos diferentes de
substituição e comparar qual deles gera menos faltas, variando também o tamanho
da memória. Este documento registra como entendemos o problema e como
pretendemos organizar o programa. Alguns pontos ainda vão ser ajustados depois
das aulas sobre o assunto.

\subsection{Entrada e saída}
\label{sec:entrada-saida}

\begin{itemize}
    \item \textbf{Entrada:} um arquivo com a sequência de páginas acessadas
    (ex.: \texttt{7 0 1 2 0 3 0 4 2 3 0 3 2 1 2 0 1 7 0 1}) e o número de
    quadros da memória física.
    \item \textbf{Saída:} o total de faltas de página, o traço mostrando o
    estado dos quadros a cada acesso e um gráfico de faltas por número de
    quadros.
\end{itemize}

% =====================================================================
\section{Conceitos que precisamos usar}
\label{sec:conceitos}

A memória virtual é dividida em blocos de tamanho fixo chamados páginas, e a
memória física é dividida em blocos do mesmo tamanho chamados quadros
\cite{silberschatz2018}. Quando o programa acessa uma página que não está em
nenhum quadro, acontece uma falta de página. Se ainda houver quadro livre, a
página só é carregada. Se a memória estiver cheia, é preciso escolher uma página
para sair, e é aí que entram os algoritmos \cite{tanenbaum2015}, resumidos na
Tabela~\ref{tab:algoritmos}.

\begin{table}[!ht]
\centering
\caption{Os cinco algoritmos de substituição que vamos implementar.}
\label{tab:algoritmos}
\small
\begin{tabular}{p{3.0cm}p{10.2cm}}
\hline
\textbf{Algoritmo} & \textbf{Qual página ele escolhe para sair} \\
\hline
FIFO & A que entrou na memória há mais tempo, sem olhar se ela está sendo usada \\
LRU & A que ficou mais tempo sem ser acessada \\
Ótimo & A que vai demorar mais para ser usada de novo (precisa conhecer o futuro) \\
Segunda Chance & Como o FIFO, mas dá uma chance extra para páginas acessadas recentemente \\
LFU & A que foi acessada menos vezes \\
\hline
\end{tabular}
\end{table}

\textbf{Anomalia de Belády:} o esperado é que aumentar a memória diminua as
faltas de página. Existe um caso em que isso não acontece e o FIFO piora com
mais quadros \cite{belady1969}. Um dos objetivos do trabalho é reproduzir esse
caso e mostrar que ele não aparece no LRU nem no Ótimo \cite{arpaci2018}.

% =====================================================================
\section{Organização dos arquivos}
\label{sec:organizacao}

Cada algoritmo fica em um arquivo separado, e todos usam a mesma assinatura de
função. Assim, para trocar de algoritmo ou mudar o número de quadros, não
precisamos mexer no resto do programa.

{\scriptsize
\begin{verbatim}
paging-simulator/
|-- src/
|   |-- main.c            /* le os parametros e chama o algoritmo   */
|   |-- simulador.h       /* tipos e assinaturas usados por todos   */
|   |-- sequencia.c       /* le o arquivo de entrada                */
|   |-- traco.c           /* imprime o traco e o resumo             */
|   |-- fifo.c            /* algoritmo FIFO                         */
|   |-- lru.c             /* algoritmo LRU                          */
|   |-- optimal.c         /* algoritmo Otimo                        */
|   |-- second_chance.c   /* algoritmo Segunda Chance               */
|   `-- lfu.c             /* algoritmo LFU                          */
|-- entradas/             /* sequencias de teste                    */
|-- analise.py            /* gera os graficos a partir dos dados    */
`-- RELATORIO.pdf
\end{verbatim}
}

% =====================================================================
\section{Funções e assinaturas}
\label{sec:funcoes}

\subsection{Tipos usados por todos os arquivos}
\label{sec:tipos}

A primeira estrutura guarda a sequência de páginas lida do arquivo de entrada:

{\scriptsize
\begin{verbatim}
typedef struct {
    int *paginas;   /* paginas na ordem em que foram acessadas */
    int  n;         /* quantas paginas a sequencia tem         */
} Sequencia;
\end{verbatim}
}

A segunda guarda o que uma simulação produziu, para um algoritmo e um número de
quadros:

{\scriptsize
\begin{verbatim}
typedef struct {
    char nome_algoritmo[32];
    int  n_frames;
    int  page_faults;
    int  hits;
} Resultado;
\end{verbatim}
}

Como todos os algoritmos devolvem esse mesmo tipo, conseguimos juntar os
resultados de todos eles em uma tabela só, sem tratar cada um de um jeito
diferente.

\subsection{Leitura da entrada}
\label{sec:leitura}

{\scriptsize
\begin{verbatim}
int  ler_sequencia(const char *caminho, Sequencia *seq);
void liberar_sequencia(Sequencia *seq);
\end{verbatim}
}

A \texttt{ler\_sequencia} abre o arquivo e guarda as páginas em um vetor de
inteiros, aceitando os números separados por espaço ou por quebra de linha;
retorna 0 se deu certo. A \texttt{liberar\_sequencia} libera esse vetor no fim
do programa.

\subsection{Impressão do traço}
\label{sec:traco}

{\scriptsize
\begin{verbatim}
void traco_linha(FILE *saida, int pagina, const int *frames,
                 int n_frames, int houve_falta, int vitima);
void traco_resumo(FILE *saida, const Resultado *r);
\end{verbatim}
}

A \texttt{traco\_linha} imprime uma linha do traço a cada acesso: qual página foi
pedida, como ficaram os quadros e se houve falta. A \texttt{traco\_resumo}
imprime o total de faltas no final. Deixamos a impressão separada dos algoritmos
para não repetir esse código cinco vezes.

\subsection{Busca nos quadros}
\label{sec:busca}

{\scriptsize
\begin{verbatim}
int buscar_frame(const int *frames, int n_frames, int pagina);
\end{verbatim}
}

Procura a página nos quadros e devolve a posição dela, ou $-1$ se ela não
estiver na memória. É essa função que diz se houve acerto ou falta, e por isso
ela é usada pelos cinco algoritmos.

\subsection{Os cinco algoritmos}
\label{sec:algoritmos}

Todos recebem os mesmos parâmetros e devolvem o mesmo tipo, mudando só a regra de
escolha da página que sai:

{\scriptsize
\begin{verbatim}
Resultado simular_fifo(const Sequencia *seq, int n_frames, FILE *traco);
Resultado simular_lru(const Sequencia *seq, int n_frames, FILE *traco);
Resultado simular_otimo(const Sequencia *seq, int n_frames, FILE *traco);
Resultado simular_segunda_chance(const Sequencia *seq, int n_frames, FILE *traco);
Resultado simular_lfu(const Sequencia *seq, int n_frames, FILE *traco);
\end{verbatim}
}

\begin{itemize}
    \item \textbf{FIFO:} pensamos em guardar um índice que aponta para o próximo
    quadro a ser substituído e ir avançando em círculo. Quando a página já está
    na memória, nada muda: a ordem de saída continua a mesma.
    \item \textbf{LRU:} a ideia é ter um contador que aumenta a cada acesso e
    anotar, para cada quadro, o valor do contador no último uso dele. Na hora de
    substituir, sai o quadro com o menor valor anotado. Assim não precisamos
    varrer toda a sequência de novo a cada acesso.
    \item \textbf{Ótimo:} como já temos a sequência inteira no vetor, dá para
    olhar para frente e ver quando cada página da memória será usada de novo.
    Sai a que demora mais, ou uma que não aparece mais. Vamos usar uma função
    auxiliar \texttt{proximo\_uso(seq, posicao, pagina)} para isso.
    \item \textbf{Segunda Chance:} além dos quadros, guardamos um bit para cada
    página, que liga quando ela é acessada. Na substituição, se o bit da página
    mais antiga estiver ligado, desligamos o bit e passamos para a próxima em vez
    de tirá-la. Se todos os bits estiverem ligados, acaba funcionando igual ao
    FIFO.
    \item \textbf{LFU:} guardamos quantas vezes cada página foi acessada e
    tiramos a de menor contagem. Em caso de empate, escolhemos a mais antiga,
    seguindo a regra do FIFO.
\end{itemize}

\subsection{Programa principal}
\label{sec:main}

{\scriptsize
\begin{verbatim}
int main(int argc, char *argv[]);
\end{verbatim}
}

Lê os parâmetros da linha de comando, carrega a sequência, escolhe qual função de
simulação chamar e imprime o resultado. A escolha do algoritmo fica em um único
ponto do código, para não espalhar condicionais pelo programa.

\subsection{Script de análise}
\label{sec:analise}

{\scriptsize
\begin{verbatim}
def plotar_faltas_por_frames(dados, saida)
def detectar_anomalia_belady(dados)
\end{verbatim}
}

A \texttt{plotar\_faltas\_por\_frames} gera o gráfico com o número de quadros no
eixo $x$ e as faltas no eixo $y$, com uma linha por algoritmo. A
\texttt{detectar\_anomalia\_belady} compara os resultados de quadros
consecutivos e avisa quando as faltas aumentaram, que é o caso da anomalia.

% =====================================================================
\section{Como o programa vai se comportar}
\label{sec:comportamento}

A execução vai ser por linha de comando, informando o arquivo de entrada, o
número de quadros e o algoritmo, e a saída esperada é a seguinte:

{\scriptsize
\begin{verbatim}
./paging-simulator --entrada entradas/sequencia1.txt --frames 3 --algoritmo fifo

Ref: 7  Frames: [7,-,-]   FAULT
Ref: 0  Frames: [7,0,-]   FAULT
Ref: 1  Frames: [7,0,1]   FAULT
Ref: 2  Frames: [2,0,1]   FAULT (substitui 7 - FIFO)
Ref: 0  Frames: [2,0,1]   HIT
...
Algoritmo: FIFO | Frames: 3 | Referencias: 20 | Faltas: 15
\end{verbatim}
}

Também vai existir uma opção para rodar todos os algoritmos de uma vez, com o
número de quadros variando de 1 até um valor escolhido. Nesse modo o traço não é
impresso; os resultados vão para um arquivo CSV que o script em Python usa para
montar o gráfico.

% =====================================================================
\section{Testes que pretendemos fazer}
\label{sec:testes}

\begin{itemize}
    \item Rodar a sequência de exemplo do enunciado com 3 e 4 quadros e conferir
    o resultado à mão, para ter certeza de que a contagem está certa.
    \item Rodar a sequência \texttt{1 2 3 4 1 2 5 1 2 3 4 5} no FIFO com 3 e
    depois com 4 quadros, que é onde esperamos ver a Anomalia de Belády.
    \item Rodar a mesma sequência no LRU e no Ótimo para mostrar que neles a
    anomalia não aparece.
    \item Gerar uma sequência maior, com páginas repetindo de tempos em tempos,
    para comparar o Segunda Chance com o FIFO puro.
    \item Variar o número de quadros de 1 até 30 em todos os algoritmos e montar
    o gráfico final.
\end{itemize}

% =====================================================================
\section{Divisão do trabalho}
\label{sec:divisao}

\begin{itemize}
    \item \textbf{Etapa 1:} estrutura dos arquivos, leitura da entrada e
    impressão do traço.
    \item \textbf{Etapa 2:} FIFO e LRU, com os testes de conferência manual.
    \item \textbf{Etapa 3:} Ótimo, Segunda Chance e LFU.
    \item \textbf{Etapa 4:} script em Python, gráficos e escrita do relatório.
\end{itemize}

% =====================================================================
\bibliographystyle{sbc}
\bibliography{Referencias/referencias}

\end{document}
